% Ragged-right paragraph column (avoids hyphenation pressure)
\newcolumntype{P}[1]{>{\raggedright\arraybackslash}p{#1}}
 
% Header style: bold, top-left, no extra gap
\renewcommand\theadfont{\bfseries\small}
\renewcommand\theadgape{}

% ==============================================================
% TABLE I — Sensors
% ==============================================================
\begin{table*}[p]
  \centering
  \footnotesize
  \renewcommand{\arraystretch}{1.4}
  \setlength{\tabcolsep}{5pt}
  \caption{Sensors chosen for the control system.}
  \label{tab:sensors}
  \begin{tabular}{|P{1.8cm}|P{2.6cm}|P{5.8cm}|P{5.8cm}|}
    \hline
    \makecell[tl]{Sensor}
      & \makecell[tl]{Item Name}
      & \makecell[tl]{Role}
      & \makecell[tl]{Implementation Details} \\
    \hline
    Water Level Sensor
      & \textbf{Robot Contact Water/Liquid Level Sensor}
      & Detects if water in the tank reaches a high or low point.
        Prevents overfilling or running dry.
      & Optical sensor output (HIGH/LOW). Connect to GPIO input.
        Can trigger pump ON/OFF automatically. \\
    \hline
    Flow Rate Sensor
      & \textbf{1/4'' Water Flow Sensor (Hall-Effect)}
      & Measures water flow rate in L/min to ensure the pump is
        moving water and to calculate total volume.
      & Generates digital pulses proportional to flow.
        Connect to GPIO input; count pulses using interrupt. \\
    \hline
    Temperature Sensor
      & \textbf{NTC Thermistor 10\,k$\Omega$ MF52-103}
      & Measures the water temperature.
        Can detect overheating or environmental changes.
      & Analog sensor. Use an ADC module (e.g.\ MCP3008) to read
        voltage, then convert to temperature in code. \\
    \hline
  \end{tabular}
\end{table*}
 
 
% ==============================================================
% TABLE II — Actuators
% ==============================================================
\begin{table*}[p]
  \centering
  \footnotesize
  \renewcommand{\arraystretch}{1.4}
  \setlength{\tabcolsep}{5pt}
  \caption{Actuators used in the control system.}
  \label{tab:actuators}
  \begin{tabular}{|P{1.8cm}|P{3.2cm}|P{5.2cm}|P{5.8cm}|}
    \hline
    \makecell[tl]{Actuator}
      & \makecell[tl]{Item Name}
      & \makecell[tl]{Role}
      & \makecell[tl]{Implementation Details} \\
    \hline
    Water Pump
      & \textbf{HiLetgo DC 12\,V 240\,L/h Micro Brushless Pump}
      & Pumps water into or out of the tank.
        Controlled automatically based on level sensor input.
      & Connect through the Pump Controller (relay/MOSFET)
        that switches the 12\,V power. \\
    \hline
    Hydraulic Valve
      & \textbf{1/4'' DC 12\,V 2-Way Normally Closed Solenoid
               Valve (WIC Valve 2ACK Series)}
      & Opens/closes to allow or stop water flow.
      & Controlled through the Valve Controller.
        When the coil is energized, the valve opens;
        otherwise, it closes. \\
    \hline
  \end{tabular}
\end{table*}
 
 
% ==============================================================
% TABLE III — Controllers
% ==============================================================
\begin{table*}[p]
  \centering
  \footnotesize
  \renewcommand{\arraystretch}{1.4}
  \setlength{\tabcolsep}{5pt}
  \caption{Modules interfacing the Raspberry Pi with actuators.}
  \label{tab:controllers}
  \begin{tabular}{|P{2cm}|P{2.2cm}|P{5.8cm}|P{5.8cm}|}
    \hline
    \makecell[tl]{Controller}
      & \makecell[tl]{Connected To}
      & \makecell[tl]{Role}
      & \makecell[tl]{Implementation} \\
    \hline
    Pump Controller
      & Water Pump
      & Acts as a switch driver---receives a 3.3\,V signal from
        the Pi to toggle 12\,V power to the pump.
      & Can be a relay board or a MOSFET circuit.
        Isolates the Pi from the high-power pump. \\
    \hline
    Valve Controller
      & Hydraulic Valve
      & Same principle as the Pump Controller, but for
        opening/closing the valve.
      & When Pi output is HIGH $\rightarrow$ coil energized
        $\rightarrow$ valve opens.
        When LOW $\rightarrow$ valve closes. \\
    \hline
  \end{tabular}
\end{table*}
 
 
% ==============================================================
% TABLE IV — SCADA / HMI
% ==============================================================
\begin{table*}[p]
  \centering
  \footnotesize
  \renewcommand{\arraystretch}{1.4}
  \setlength{\tabcolsep}{5pt}
  \caption{SCADA / Human--Machine Interface components.}
  \label{tab:scada}
  \begin{tabular}{|P{2.4cm}|P{2.4cm}|P{11cm}|}
    \hline
    \makecell[tl]{Component}
      & \makecell[tl]{Role}
      & \makecell[tl]{Description} \\
    \hline
    Laptop / Computer
      & CLI or GUI
      & Acts as the human--machine interface (HMI) for remote
        monitoring and manual control. \\
    \hline
    Wireless Communication
      & Local Area Wifi
      & Connects the Raspberry Pi to the laptop using Wi-Fi.
        Can send real-time data or receive control commands. \\
    \hline
    Software Options
      & Data Visualization
      & \textbf{Flask Web Server:} Allows CLI or GUI to send commands to SCADA.
        \textbf{Plotly:} Allows GUI to poll SQL database to generate plots.
        \textbf{SSH:} Remote access for starting and troubleshooting server.\\
    \hline
  \end{tabular}
\end{table*}
 
 
% ==============================================================
% TABLE V — Communication Links (Pi--Arduino)
% ==============================================================
\begin{table*}[p]
  \centering
  \footnotesize
  \renewcommand{\arraystretch}{1.4}
  \setlength{\tabcolsep}{5pt}
  \caption{Communication link options between the Raspberry Pi and Arduino.}
  \label{tab:comms}
  \begin{tabular}{|P{3.2cm}|P{3cm}|P{3cm}|P{6.6cm}|}
    \hline
    \makecell[tl]{Link}
      & \makecell[tl]{Pros}
      & \makecell[tl]{Cons}
      & \makecell[tl]{Method} \\
    \hline
    \textbf{USB Serial (UART over USB)}
    \textit{(Recommended)}
      & Easiest; powered and isolated via cable; stable.
      & Requires a USB cable.
      & Connect Arduino to Pi via USB; use
        \texttt{/dev/ttyACM0} or \texttt{/dev/ttyUSB0}
        at 115200\,baud. \\
    \hline
    \textbf{I$^{2}$C} (Pi master, Arduino slave)
      & Simple 2-wire interface.
      & 3.3$\leftrightarrow$5\,V level shifting; address management.
      & Pi SDA/SCL $\leftrightarrow$ Arduino A4/A5 (Uno);
        add bi-directional level shifter. \\
    \hline
    \textbf{SPI}
      & Fast throughput.
      & Chip-select management; level shifting required.
      & Pi SPI0 $\leftrightarrow$ Arduino SPI pins;
        level shifter required. \\
    \hline
    \textbf{RS-485 (Modbus RTU)}
      & Long cable runs; noise-robust.
      & Requires transceivers and protocol stack.
      & MAX485 modules, twisted-pair cable, Modbus library. \\
    \hline
  \end{tabular}
\end{table*}